\documentclass[11pt,a4paper]{article}

\usepackage[utf8]{inputenc}
\usepackage[T1]{fontenc}
\usepackage[french]{babel}
\usepackage[margin=2.5cm]{geometry}
\usepackage{xcolor}
\usepackage{listings}
\usepackage{amsmath}
\usepackage{booktabs}
\usepackage{array}
\usepackage{tikz}
\usepackage{hyperref}
\usepackage{enumitem}
\usepackage{parskip}

\definecolor{codebg}{RGB}{245,245,245}
\definecolor{codekw}{RGB}{0,90,150}
\definecolor{codestr}{RGB}{150,60,20}
\definecolor{codecom}{RGB}{110,110,110}

\lstset{
  backgroundcolor=\color{codebg},
  basicstyle=\ttfamily\small,
  keywordstyle=\color{codekw}\bfseries,
  stringstyle=\color{codestr},
  commentstyle=\color{codecom}\itshape,
  breaklines=true,
  showstringspaces=false,
  frame=single,
  rulecolor=\color{gray!40},
  language=Python,
  xleftmargin=1em,
  xrightmargin=1em,
  aboveskip=1em,
  belowskip=1em,
}

\hypersetup{colorlinks=true, linkcolor=blue!50!black, urlcolor=blue!50!black}

\title{\Large Reconnaissance d'image par CNN\\[0.3em]\normalsize Jour 1 --- Images, tenseurs et donn\'ees}
\author{}
\date{}

\begin{document}
\maketitle
\vspace{-2em}
\tableofcontents

\section{Ce qu'on fait aujourd'hui et pourquoi}

Avant d'entra\^iner quoi que ce soit, il faut comprendre ce qu'on donne
\`a manger \`a un r\'eseau de neurones. Un CNN (r\'eseau de neurones
convolutif) ne \guillemotleft~voit~\guillemotright\ pas une image comme nous : il manipule des
tableaux de nombres. Le Jour~1 sert uniquement \`a construire cette
intuition et \`a pr\'eparer les donn\'ees, \emph{avant} de toucher au
mod\`ele (Jour~2).

Le fichier concern\'e est \texttt{src/data.py}. \`A la fin de cette
lecture, tu devrais savoir r\'epondre \`a~:
\begin{itemize}[nosep]
  \item Comment une image est stock\'ee en m\'emoire (pixels, canaux) ?
  \item Qu'est-ce qu'un \emph{tenseur} et pourquoi ce mot revient partout ?
  \item Pourquoi on divise les pixels par 255 ?
  \item Pourquoi on s\'epare les donn\'ees en \guillemotleft~train~\guillemotright\ et \guillemotleft~test~\guillemotright\ ?
  \item Que contient exactement CIFAR-10 ?
\end{itemize}

\section{Comment un ordinateur repr\'esente une image}

\subsection{Le pixel}

Une image num\'erique est une grille de petits carr\'es de couleur
appel\'es \emph{pixels} (\emph{picture element}). Une image de
32~$\times$~32 contient $32 \times 32 = 1024$ pixels, dispos\'es en
32~lignes et 32~colonnes.

\subsection{La couleur : le mod\`ele RGB}

Chaque pixel couleur est cod\'e par trois nombres entre 0 et 255~:
\begin{itemize}[nosep]
  \item \textbf{R} (Red) --- intensit\'e de rouge
  \item \textbf{G} (Green) --- intensit\'e de vert
  \item \textbf{B} (Blue) --- intensit\'e de bleu
\end{itemize}

En m\'elangeant ces trois intensit\'es on obtient n'importe quelle
couleur. Par exemple $(255, 0, 0)$ est un rouge pur, $(0, 0, 0)$ est
noir, $(255, 255, 255)$ est blanc, $(255, 165, 0)$ est orange.

\begin{center}
\begin{tikzpicture}[x=1cm,y=1cm]
  \foreach \i in {0,...,4} {
    \foreach \j in {0,...,3} {
      \pgfmathsetmacro{\r}{50+\i*40}
      \pgfmathsetmacro{\g}{80+\j*40}
      \definecolor{tmpcol}{RGB}{\r,\g,180}
      \fill[tmpcol] (\i,\j) rectangle ++(1,1);
      \draw[white,line width=0.5pt] (\i,\j) rectangle ++(1,1);
    }
  }
\end{tikzpicture}
\\[0.3em]
{\small Une image n'est qu'une grille de pixels, chacun \'etant un triplet de nombres.}
\end{center}

Une image couleur de taille $H \times L$ (hauteur $\times$ largeur) est
donc repr\'esent\'ee par $H \times L \times 3$ nombres --- le \guillemotleft~3~\guillemotright\
correspond aux trois \emph{canaux} R, G, B. Pour CIFAR-10, chaque image
fait $32 \times 32 \times 3 = 3072$ nombres.

\section{Le tenseur : la structure de donn\'ees du deep learning}

Le mot \emph{tenseur} revient partout en deep learning (le \guillemotleft~T~\guillemotright\ de
Tensor\textbf{Flow}). C'est simplement une g\'en\'eralisation de notions
que tu connais d\'ej\`a~:

\begin{center}
\begin{tabular}{@{}lll@{}}
\toprule
Dimension & Nom usuel & Exemple \\
\midrule
0 & scalaire & \texttt{5} \\
1 & vecteur & \texttt{[5, 3, 1]} \\
2 & matrice & tableau de nombres (lignes/colonnes) \\
3+ & tenseur & empilement de matrices \\
\bottomrule
\end{tabular}
\end{center}

En pratique, en Python/NumPy/TensorFlow, tout --- scalaire, vecteur,
matrice ou tenseur --- est repr\'esent\'e par le m\^eme objet
(\texttt{np.ndarray} ou \texttt{tf.Tensor}), caract\'eris\'e par sa
\textbf{forme} (\emph{shape}) : la liste des tailles de chaque dimension.

Une image couleur est donc un tenseur de forme $(32, 32, 3)$ : 32 lignes,
32 colonnes, 3 canaux. Et un \emph{lot} (\emph{batch}) de 50~000 images
d'entra\^inement est un tenseur de forme $(50000, 32, 32, 3)$ : c'est
exactement ce que renvoie \texttt{cifar10.load\_data()}.

\textbf{Pourquoi cette dimension suppl\'ementaire pour le batch ?} Parce
qu'un r\'eseau de neurones ne traite jamais une image \`a la fois : il
en traite plusieurs en parall\`ele (par exemple 64), ce qui est bien plus
rapide sur CPU/GPU gr\^ace au calcul matriciel. Cette dimension s'appelle
souvent \texttt{N} (nombre d'exemples) dans la convention
\texttt{(N, H, W, C)}.

\section{Le dataset CIFAR-10}

CIFAR-10 est un jeu de donn\'ees standard, cr\'e\'e \`a l'universit\'e de
Toronto, tr\`es utilis\'e pour apprendre la classification d'images
parce qu'il est petit, rapide \`a entra\^iner et d\'ej\`a \'etiquet\'e.

\begin{center}
\begin{tabular}{@{}ll@{}}
\toprule
Caract\'eristique & Valeur \\
\midrule
Nombre total d'images & 60~000 \\
Images d'entra\^inement (\emph{train}) & 50~000 \\
Images de test (\emph{test}) & 10~000 \\
Taille de chaque image & $32 \times 32$ pixels, RGB \\
Nombre de classes & 10 \\
\bottomrule
\end{tabular}
\end{center}

Les 10 classes, avec leur indice num\'erique (c'est cet indice qui est
r\'eellement stock\'e, pas le mot) :

\begin{center}
\begin{tabular}{@{}cl@{}c@{}cl@{}}
\toprule
Indice & Classe & \phantom{xx} & Indice & Classe \\
\midrule
0 & avion & & 5 & chien \\
1 & voiture & & 6 & grenouille \\
2 & oiseau & & 7 & cheval \\
3 & chat & & 8 & bateau \\
4 & cerf & & 9 & camion \\
\bottomrule
\end{tabular}
\end{center}

C'est exactement la liste \texttt{CLASS\_NAMES} dans \texttt{src/data.py}.

\section{Les labels : pourquoi de simples entiers}

Pour chaque image, CIFAR-10 fournit un \emph{label} (\'etiquette) : un
entier entre 0 et 9 indiquant la bonne classe. Il existe deux fa\c{c}ons
courantes d'encoder \c{c}a en deep learning~:

\begin{itemize}
  \item \textbf{Encodage entier} (\emph{sparse}) : le label est
  directement \texttt{3} pour \guillemotleft~chat~\guillemotright.
  \item \textbf{One-hot} : le label est un vecteur de 10 valeurs, toutes
  \`a 0 sauf un 1 \`a la position de la bonne classe --- pour \guillemotleft~chat~\guillemotright\
  (indice 3) : \texttt{[0,0,0,1,0,0,0,0,0,0]}.
\end{itemize}

Les deux repr\'esentent la m\^eme information. On garde ici l'encodage
entier (plus simple, moins de m\'emoire), et c'est pour \c{c}a que
\texttt{model.py} utilise la fonction de perte
\texttt{sparse\_categorical\_\allowbreak crossentropy} plut\^ot que
\texttt{categorical\_crossentropy} (qui, elle, attendrait du one-hot). On
reverra cette histoire de \guillemotleft~fonction de perte~\guillemotright\ en d\'etail au
Jour~2.

\section{La normalisation : pourquoi diviser par 255}

Dans \texttt{data.py}~:
\begin{lstlisting}
x_train = x_train.astype("float32") / 255.0
\end{lstlisting}

Chaque pixel est \`a l'origine un entier entre 0 et 255. On le convertit
en nombre \`a virgule (\texttt{float32}) puis on le ram\`ene entre 0 et 1
en divisant par 255 (la valeur maximale possible).

\textbf{Pourquoi c'est important ?} Un r\'eseau de neurones apprend en
ajustant progressivement ses param\`etres via un algorithme appel\'e
\emph{descente de gradient} (on d\'etaillera au Jour~2). Cet algorithme
est beaucoup plus stable et converge plus vite quand toutes les entr\'ees
sont \`a une \'echelle comparable, proche de 0. Avec des valeurs brutes
allant jusqu'\`a 255, les calculs internes (multiplications, sommes)
produisent des nombres bien plus grands, ce qui peut :

\begin{itemize}[nosep]
  \item ralentir l'apprentissage (il faut plus d'it\'erations pour
  converger) ;
  \item le rendre instable (les mises \`a jour des param\`etres peuvent
  \guillemotleft~sauter~\guillemotright\ au lieu de converger doucement) ;
  \item saturer certaines fonctions d'activation qui \'ecrasent les
  grandes valeurs (moins de signal utile pour apprendre).
\end{itemize}

C'est une des r\`egles les plus universelles du deep learning : \textbf{on
normalise quasiment toujours les entr\'ees}. Pour des images, diviser par
255 est la m\'ethode la plus simple et suffit largement ici.

\section{Pourquoi s\'eparer donn\'ees d'entra\^inement et de test}

\texttt{cifar10.load\_data()} renvoie directement deux ensembles
s\'epar\'es. C'est volontaire et fondamental~:

\begin{itemize}
  \item Le mod\`ele apprend uniquement sur les 50~000 images de
  \textbf{train}.
  \item On \'evalue sa performance sur les 10~000 images de
  \textbf{test}, qu'il n'a \emph{jamais vues} pendant l'entra\^inement.
\end{itemize}

Pourquoi ? Un mod\`ele suffisamment complexe peut \guillemotleft~apprendre par
c\oe{}ur~\guillemotright\ les images d'entra\^inement (il m\'emorise plut\^ot qu'il ne
g\'en\'eralise). Un tel mod\`ele obtiendrait un score parfait sur les
donn\'ees qu'il a d\'ej\`a vues, mais serait mauvais sur une image
nouvelle --- inutile en pratique. Ce ph\'enom\`ene s'appelle
l'\textbf{overfitting} (surapprentissage) ; on le reverra concr\`etement
au Jour~2 quand on comparera l'accuracy sur train et sur test pendant
l'entra\^inement.

\`A retenir : \emph{la seule mesure de performance qui compte, c'est
celle sur des donn\'ees jamais vues}.

\section{Lecture comment\'ee de \texttt{src/data.py}}

\begin{lstlisting}
import tensorflow as tf

CLASS_NAMES = [
    "avion", "voiture", "oiseau", "chat", "cerf",
    "chien", "grenouille", "cheval", "bateau", "camion",
]

def load_data():
    (x_train, y_train), (x_test, y_test) = tf.keras.datasets.cifar10.load_data()

    # normalisation : 0-255 -> 0-1 (voir section 6)
    x_train = x_train.astype("float32") / 255.0
    x_test = x_test.astype("float32") / 255.0

    # y_train/y_test arrivent en colonnes (N, 1), on les aplatit en (N,)
    y_train = y_train.flatten()
    y_test = y_test.flatten()

    return (x_train, y_train), (x_test, y_test)
\end{lstlisting}

D\'etail du seul point pas encore expliqu\'e : \texttt{flatten()}. Keras
renvoie les labels sous forme de tableau \`a deux dimensions, de forme
$(50000, 1)$ --- une colonne. \texttt{flatten()} le transforme en tableau
\`a une dimension, de forme $(50000,)$, ce qui est plus simple \`a
manipuler (par exemple \texttt{y\_train[0]} donne directement l'entier du
label, pas un tableau \`a un \'el\'ement).

\section{\`A faire toi-m\^eme}

Sur ta machine (l'environnement d'ex\'ecution distant n'a pas acc\`es au
serveur qui h\'eberge CIFAR-10) :

\begin{enumerate}
  \item Lance :
\begin{lstlisting}
cd src
python data.py
\end{lstlisting}
  Tu dois voir s'afficher la forme des tableaux
  (\texttt{Train: (50000, 32, 32, 3)}, \texttt{Test: (10000, 32, 32, 3)})
  et un exemple de label.

  \item Visualise quelques images pour v\'erifier que tout correspond
  bien, avec ce petit script \`a coller dans un fichier \`a part
  (par exemple \texttt{src/explore.py}) :
\begin{lstlisting}
import matplotlib.pyplot as plt
from data import load_data, CLASS_NAMES

(x_train, y_train), _ = load_data()

fig, axes = plt.subplots(2, 5, figsize=(10, 4))
for i, ax in enumerate(axes.flat):
    ax.imshow(x_train[i])
    ax.set_title(CLASS_NAMES[y_train[i]])
    ax.axis("off")
plt.tight_layout()
plt.show()
\end{lstlisting}
  Tu devrais voir 10 petites images 32$\times$32 avec leur nom de classe
  au-dessus. C'est la meilleure fa\c{c}on de \guillemotleft~voir~\guillemotright\ concr\`etement
  ce que le mod\`ele va recevoir en entr\'ee.
\end{enumerate}

\section{R\'ecapitulatif}

\begin{center}
\begin{tabular}{@{}p{4cm}p{10cm}@{}}
\toprule
Concept & En une phrase \\
\midrule
Pixel & Un point de l'image, cod\'e par 3 nombres (R, G, B) entre 0 et 255. \\
Tenseur & Un tableau de nombres \`a $n$ dimensions, caract\'eris\'e par sa forme (\emph{shape}). \\
Shape $(N,H,W,C)$ & $N$ images, hauteur $H$, largeur $W$, $C$ canaux de couleur. \\
Normalisation & Ramener les pixels entre 0 et 1 pour un apprentissage plus rapide et stable. \\
Label entier & L'indice (0-9) de la bonne classe, sans encodage one-hot. \\
Train / test & S\'eparer les donn\'ees pour mesurer si le mod\`ele g\'en\'eralise, pas seulement m\'emorise. \\
\bottomrule
\end{tabular}
\end{center}

\section{La suite : Jour 2}

Une fois ces notions solides, on passera \`a \texttt{src/model.py} et
\texttt{src/train.py} : construire le CNN lui-m\^eme (convolution,
filtre, pooling), lancer un premier entra\^inement, et comprendre ce que
veulent dire \emph{accuracy}, \emph{loss}, \emph{overfitting} et
\emph{dropout} en pratique, sur des courbes r\'eelles.

\end{document}
